\documentclass[10pt,a4paper]{amsart}
\usepackage{amsmath,amssymb,amsthm,graphicx}
\usepackage{verbatim}
\newenvironment{reference}{\comment}{\endcomment}
%\newenvironment{reference}{}{}
\newenvironment{slogan}{\comment}{\endcomment}
\newenvironment{history}{\comment}{\endcomment}

% For commutative diagrams we use Xy-pic
\usepackage[all]{xy}

% We use 2cell for 2-commutative diagrams.
\xyoption{2cell}
% Theorem environments.
%
\theoremstyle{plain}
\newtheorem{theorem}[subsection]{Theorem}
\newtheorem{proposition}[subsection]{Proposition}
\newtheorem{lemma}[subsection]{Lemma}

\theoremstyle{definition}
\newtheorem{definition}[subsection]{Definition}
\newtheorem{example}[subsection]{Example}
\newtheorem{exercise}[subsection]{Exercise}
\newtheorem{situation}[subsection]{Situation}

\theoremstyle{remark}
\newtheorem{remark}[subsection]{Remark}
\newtheorem{remarks}[subsection]{Remarks}

\numberwithin{equation}{subsection}

% Macros
%
\def\lim{\mathop{\mathrm{lim}}\nolimits}
\def\colim{\mathop{\mathrm{colim}}\nolimits}
\def\Spec{\mathop{\mathrm{Spec}}}
\def\Hom{\mathop{\mathrm{Hom}}\nolimits}
\def\Ext{\mathop{\mathrm{Ext}}\nolimits}
\def\SheafHom{\mathop{\mathcal{H}\!\mathit{om}}\nolimits}
\def\SheafExt{\mathop{\mathcal{E}\!\mathit{xt}}\nolimits}
\def\Sch{\mathit{Sch}}
\def\Mor{\mathop{\mathrm{Mor}}\nolimits}
\def\Ob{\mathop{\mathrm{Ob}}\nolimits}
\def\Sh{\mathop{\mathit{Sh}}\nolimits}
\def\NL{\mathop{N\!L}\nolimits}
\def\CH{\mathop{\mathrm{CH}}\nolimits}
\def\proetale{{pro\text{-}\acute{e}tale}}
\def\etale{{\acute{e}tale}}
\def\QCoh{\mathit{QCoh}}
\def\Ker{\mathop{\mathrm{Ker}}}
\def\Im{\mathop{\mathrm{Im}}}
\def\Coker{\mathop{\mathrm{Coker}}}
\def\Coim{\mathop{\mathrm{Coim}}}

% Boxtimes
%
\DeclareMathSymbol{\boxtimes}{\mathbin}{AMSa}{"02}

%
% Macros for moduli stacks/spaces
%
\def\QCohstack{\mathcal{QC}\!\mathit{oh}}
\def\Cohstack{\mathcal{C}\!\mathit{oh}}
\def\Spacesstack{\mathcal{S}\!\mathit{paces}}
\def\Quotfunctor{\mathrm{Quot}}
\def\Hilbfunctor{\mathrm{Hilb}}
\def\Curvesstack{\mathcal{C}\!\mathit{urves}}
\def\Polarizedstack{\mathcal{P}\!\mathit{olarized}}
\def\Complexesstack{\mathcal{C}\!\mathit{omplexes}}
% \Pic is the operator that assigns to X its picard group, usage \Pic(X)
% \Picardstack_{X/B} denotes the Picard stack of X over B
% \Picardfunctor_{X/B} denotes the Picard functor of X over B
\def\Pic{\mathop{\mathrm{Pic}}\nolimits}
\def\Picardstack{\mathcal{P}\!\mathit{ic}}
\def\Picardfunctor{\mathrm{Pic}}
\def\Deformationcategory{\mathcal{D}\!\mathit{ef}}
\usepackage{cleveref}
\crefname{originallemma}{original lemma}{original lemmas}
\Crefname{originallemma}{Original lemma}{Original lemmas}
\makeatletter
\newlabel{lemma-lift-map}{{0G2H}{source}{}{sourcepointer-1}{}}
\newlabel{lemma-lift-map@cref}{{[originallemma][0][]0G2H}{[1][0][]source}{}{}{}}
\makeatother
\crefname{originalreference}{original reference}{original references}
\Crefname{originalreference}{Original reference}{Original references}
\makeatletter
\newlabel{categories-remark-pro-category-copresheaves}{{0G2X}{source}{}{sourcepointer-2}{}}
\newlabel{categories-remark-pro-category-copresheaves@cref}{{[originalreference][0][]0G2X}{[1][0][]source}{}{}{}}
\makeatother
\crefname{originallemma}{original lemma}{original lemmas}
\Crefname{originallemma}{Original lemma}{Original lemmas}
\makeatletter
\newlabel{perfect-lemma-lift-coherent}{{0G48}{source}{}{sourcepointer-3}{}}
\newlabel{perfect-lemma-lift-coherent@cref}{{[originallemma][0][]0G48}{[1][0][]source}{}{}{}}
\makeatother
\crefname{originalproposition}{original proposition}{original propositions}
\Crefname{originalproposition}{Original proposition}{Original propositions}
\makeatletter
\newlabel{perfect-proposition-DCoh}{{0FDB}{source}{}{sourcepointer-4}{}}
\newlabel{perfect-proposition-DCoh@cref}{{[originalproposition][0][]0FDB}{[1][0][]source}{}{}{}}
\makeatother
\title{Lifting a bounded coherent object from an open to a Deligne system, uniquely}
\author{The Stacks Project Authors}
\date{stacks-project:duality.tex:local-snapshot-20261012; GNU FDL 1.2 or later}
\begin{document}
\maketitle
\noindent\textit{Editorial note.} All delivered mathematical text is the author's own, copied from the pinned Stacks Project chapter file duality.tex; no statement or proof step was rewritten, reordered or removed.\par\medskip
\noindent\textit{Editorial note.} Preamble provenance: the pinned chapter file opens with an input of the shared preamble and carries no preamble text of its own, so the project's single Stacks-wide preamble staged at sources/stacks-duality/preamble.tex is excerpted verbatim under the recipe fixed by LANE\_BRIEF 22.1, as three spans, lines 8-13 (verbatim package and the reference/slogan/history environments), lines 14-18 (Xy-pic setup) and lines 154-217 (the native theorem declarations and mathematical macros). Nothing was added, removed or redefined; the wrapper class and the amsmath/amssymb/amsthm/graphicx base packages are supplied by the assembler.\par\medskip
\noindent\textit{Editorial note.} The counted statement and proof contain no cross reference, no citation and no included file, so the delivered body is self contained.\par\medskip
\noindent\textit{Editorial note.} Canonical location of the counted proposition: Stacks tag 0G4L (sources/stacks-duality/pinned-original-source.tex line 8103).\par\medskip
\noindent\textit{Allowed standard prerequisites (editorial).}\begin{itemize}
\item The colimit of Hom out of a Deligne system and the pro-category description. Same chapter, Lemma labelled lemma-lift-map, cited at proof lines 8144-8146, and Categories, Remark labelled categories-remark-pro-category-copresheaves, cited at proof lines 8146-8150. Gives part (1).
\item Lifting a bounded coherent object from an open. Derived Categories of Schemes, Lemma labelled perfect-lemma-lift-coherent, cited at proof lines 8156-8158. Produces K' restricting to K.
\item Representation of bounded coherent objects by complexes of coherent modules. Derived Categories of Schemes, Proposition labelled perfect-proposition-DCoh, cited at proof lines 8158-8160. Gives the complex whose ideal powers define the Deligne system.
\end{itemize}\medskip
\ifdefined\hypertarget\hypertarget{sourcepointer-1}{}\fi
\noindent\textit{Source reference.} Original lemma 0G2H: sources/stacks-duality/pinned-original-source.tex, source lines 7920--7920. Same-chapter pointer into the author environment lemma at source line 7920; invoked at proof line 8121 of the counted proof. It is retained only as a navigation pointer; every load-bearing step is inside the delivered span.\par\medskip
\ifdefined\hypertarget\hypertarget{sourcepointer-2}{}\fi
\noindent\textit{Source reference.} Original reference 0G2X: sources/stacks-chow/tags.txt, source lines 20072--20072. Cross-chapter standard prerequisite located in the project tag list staged in the local snapshot; invoked at proof line 8124 of the counted proof. It is retained only as a navigation pointer; every load-bearing step is inside the delivered span.\par\medskip
\ifdefined\hypertarget\hypertarget{sourcepointer-3}{}\fi
\noindent\textit{Source reference.} Original lemma 0G48: sources/stacks-chow/tags.txt, source lines 20118--20118. Cross-chapter standard prerequisite located in the project tag list staged in the local snapshot; invoked at proof line 8129 of the counted proof. It is retained only as a navigation pointer; every load-bearing step is inside the delivered span.\par\medskip
\ifdefined\hypertarget\hypertarget{sourcepointer-4}{}\fi
\noindent\textit{Source reference.} Original proposition 0FDB: sources/stacks-chow/tags.txt, source lines 19211--19211. Cross-chapter standard prerequisite located in the project tag list staged in the local snapshot; invoked at proof line 8130 of the counted proof. It is retained only as a navigation pointer; every load-bearing step is inside the delivered span.\par\medskip
\begin{lemma}
\label{lemma-extension-by-zero}
Let $j : U \to X$ be an open immersion of Noetherian schemes.
\begin{enumerate}
\item Let $(K_n)$ and $(L_n)$ be Deligne systems.
Let $K$ and $L$ be the values of the constant systems
$(K_n|_U)$ and $(L_n|_U)$. Given a morphism $\alpha : K \to L$
of $D(\mathcal{O}_U)$
there is a unique morphism of pro-systems $(K_n) \to (L_n)$
of $D^b_{\textit{Coh}}(\mathcal{O}_X)$ whose restriction to $U$ is $\alpha$.
\item Given $K \in D^b_{\textit{Coh}}(\mathcal{O}_U)$ there exists a
Deligne system $(K_n)$ such that $(K_n|_U)$ is constant
with value $K$.
\item The pro-object $(K_n)$ of $D^b_{\textit{Coh}}(\mathcal{O}_X)$
of (2) is unique up to unique isomorphism (as a pro-object).
\end{enumerate}
\end{lemma}

\begin{proof}
Part (1) is an immediate consequence of Lemma \ref{lemma-lift-map}
and the fact that morphisms between pro-systems are
the same as morphisms between the functors they corepresent, see
Categories, Remark \ref{categories-remark-pro-category-copresheaves}.

\medskip\noindent
Let $K$ be as in (2). We can choose $K' \in D^b_{\textit{Coh}}(\mathcal{O}_X)$
whose restriction to $U$ is isomorphic to $K$, see
Derived Categories of Schemes, Lemma \ref{perfect-lemma-lift-coherent}.
By Derived Categories of Schemes, Proposition \ref{perfect-proposition-DCoh}
we can represent $K'$ by a bounded complex $\mathcal{F}^\bullet$
of coherent $\mathcal{O}_X$-modules. Choose a quasi-coherent sheaf
of ideals $\mathcal{I} \subset \mathcal{O}_X$ whose vanishing
locus is $X \setminus U$ (for example choose $\mathcal{I}$ to correspond
to the reduced induced subscheme structure on $X \setminus U$).
Then we can set $K_n$ equal to the object represented by the complex
$\mathcal{I}^n\mathcal{F}^\bullet$ as in the introduction
to this section.

\medskip\noindent
Part (3) is immediate from parts (1) and (2).
\end{proof}

\end{document}
